\pagestyle{myheadings} \setcounter{page}{1} \setcounter{footnote}{0}

\section{~与 NUOPC 耦合} \label{app:nuopc}
\newcounters

\vssub
\subsection{~引言} \label{sec:nuopcintro}
\vssub

从 v6.02 起，\ws\ 具有一个分量 cap {\code wmesmf}，它通过 National
Unified Operational Prediction Capability（NUOPC）
\footnote{https://earthsystemcog.org/projects/nuopc/} Layer 与多网格例程
建立接口；NUOPC 规定使用 Earth System Modeling Framework（ESMF）
\footnote{https://www.earthsystemcog.org/projects/esmf/} 与大气、海洋、冰和
风暴潮模式等其他地球系统耦合。该 cap 旨在保持灵活，并已在 NOAA 和美国
海军的多个不同耦合模式中使用。本节说明如何构建、安装、配置和运行 NUOPC
cap。它假定读者具备 \ws\ 的基本知识。

\vssub
\subsection{~构建和安装 NUOPC Cap} \label{sec:nuopcbuild}
\vssub

若要生成一个库，使其中包含可纳入 NUOPC 耦合模式的 \ws\ 代码和 cap，请使用
{\code model\/esmf\/Makefile}。例如，在 NOAA Environmental Modeling System
（NEMS）中，命令为：
\command{make ww3\_nems}
该 makefile 随后会调用 {\code w3\_make}。作为这一过程的一部分，还会创建
一个 nuopc.mk makefile 片段，用来告诉 NUOPC/ESMF \ws\ 库的位置。

注意有一个新的开关 {\code WRST}，它会把 10 m 风加入重启动文件，并在波浪
模式初始时间步使用该风场。这可用于耦合大气模式在初始化时没有 10 m 风速的
情况。

\vssub
\subsection{~NUOPC Cap 中的导入/导出场} \label{sec:nuopcfields}
\vssub

可用的导入和导出场列在下面。有关如何为某个导入场激活耦合的信息，请参见
第~\ref{sec:nuopcconfig} 节。

\noindent 导入场：
\begin{itemize}
\item sea\_surface\_height\_above\_sea\_level
\item surface\_eastward\_sea\_water\_velocity
\item surface\_northward\_sea\_water\_velocity
\item eastward\_wind\_at\_10m\_height
\item northward\_wind\_at\_10m\_height
\item sea\_ice\_concentration
\end{itemize}

\noindent 导出：
\begin{itemize}
\item wave\_induced\_charnock\_parameter
\item wave\_z0\_roughness\_length
\item eastward\_stokes\_drift\_current
\item northward\_stokes\_drift\_current
\item eastward\_wave\_bottom\_current
\item northward\_wave\_bottom\_current
\item wave\_bottom\_current\_period
\item eastward\_wave\_radiation\_stress
\item eastward\_northward\_wave\_radiation\_stress
\item northward\_wave\_radiation\_stress
\end{itemize}


\vssub
\subsection{~NUOPC Cap 输入文件配置} \label{sec:nuopcconfig}
\vssub

所需的 \ws\ 输入文件是 ww3\_multi.inp 或 ww3\_multi.nml 文件。若要指定某个
特定输入场通过耦合而非文件获取，需要在该输入场的输入网格说明前放置
`CPL:'。

注意，NUOPC cap 当前的限制是只能有一个输入网格（无论它是模式网格之一还是
输入网格）和一个导出网格（如果有多个计算网格，则为第一个计算网格）。网格
可以是非结构网格或结构网格。

注意，虽然运行的起止时间由 NUOPC 驱动程序决定，输出的起始时间、结束时间
和频率仍由 ww3\_multi 输入文件决定。

\vssub
\subsection{~运行 NUOPC Cap} \label{sec:nuopcrun}
\vssub

该 cap 设计为用于耦合系统，这超出了本文档的范围。{\code
regtest/run\_esmf\_test\_suite} 脚本提供了一个运行独立 \ws\ cap 回归测试
的脚本。
